\chapter{La década de Raúl}
\label{ch:raul}

\section{Otra manera de gobernar}

Fidel Castro enfermó el 31 de julio de 2006 y nunca volvió al cargo; Raúl Castro
asumió provisionalmente ese mismo día y formalmente en febrero de 2008. El cambio
de estilo fue inmediato y mayor que el cambio de política.

Fidel gobernaba por discurso: los problemas se nombraban en público, las
soluciones se anunciaban como campañas, y la campaña era el instrumento. Raúl
gobernaba por reunión. Delegaba, ponía plazos, exigía que los ministerios
respondieran por sus metas, y ---lo más inusual--- decía en público que el sistema
no funcionaba. El discurso de Camagüey del 26 de julio de 2007 es el texto: le
dijo a su auditorio que los salarios cubanos no alcanzaban para vivir, que hacían
falta cambios estructurales y de concepto, y, en el ejemplo que todos recuerdan,
que un país incapaz de poner un vaso de leche en la mesa de un niño todos los días
tenía un problema que debía dejar de explicar y empezar a resolver. Viniendo del
segundo hombre de la revolución, en el discurso del aniversario, aquello fue
licencia para discutir cosas que no habían sido discutibles.

Abrió además una consulta. A finales de 2010 el borrador de los lineamientos
económicos circuló para su discusión en centros de trabajo y barrios; el gobierno
informó de unas 160.000 reuniones y millones de comentarios, y de que alrededor de
dos tercios de los párrafos del borrador se modificaron
\citep{mesalago2013}.\unv{} Sea cual sea el juicio sobre el mecanismo, fue la
consulta formal más amplia que había hecho jamás el Estado cubano, y los cubanos
participaron en serio porque creyeron que algo saldría de ella.

\section{La retirada de las humillaciones}

Las primeras medidas no fueron reformas económicas en absoluto. Fueron la
derogación de prohibiciones cuya única función había sido recordarles a los
cubanos su lugar.

En 2008 se permitió a los cubanos comprar teléfonos móviles, computadoras,
reproductores de DVD y microondas, y hospedarse en hoteles turísticos ---poniendo
fin, doce años tarde, a la situación en que un cubano no podía entrar en un
edificio en suelo cubano donde sí podía entrar un extranjero---. En 2011 se les
permitió comprar y vender sus propias casas y automóviles, prohibido desde 1960 y
que había producido un elaborado mercado gris de permutas, matrimonios ficticios y
sobornos.

Y en octubre de 2012, con efecto desde enero de 2013, se abolió el permiso de
salida. La tarjeta blanca ---el permiso para salir, que costaba dinero, tardaba
meses, exigía carta de invitación y se denegaba a discreción--- desapareció, y un
cubano podía viajar con su pasaporte como cualquiera. Los médicos y algunas otras
categorías siguieron restringidos en la práctica, y la prohibición de retorno de
ocho años se redujo a dos. Pero el principio cambió: después de cincuenta y cuatro
años, salir temporalmente de Cuba dejó de ser un acto político.

Nada de esto le costó nada al Estado. Ese es precisamente el punto, y el
capítulo~\ref{ch:sequence} lo convierte en principio operativo: la primera fase de
cualquier programa de reformas debería empezar por las prohibiciones que no cuesta
nada retirar, porque son gratis, son inmediatas, y compran la credibilidad que las
medidas caras van a necesitar. Los actos más populares de Raúl Castro fueron todos
de esa clase.

\section{Los lineamientos, y por qué se detuvieron}

El VI Congreso del Partido, de abril de 2011 ---el primero en catorce años---,
adoptó 313 Lineamientos de la Política Económica y Social. El programa que
describían era coherente: reducir la plantilla estatal, ampliar el empleo no
estatal, descentralizar las decisiones empresariales, separar la administración
estatal de la gestión empresarial, unificar la moneda y pasar del subsidio
universal a la asistencia focalizada. El primer paso anunciado era el traslado de
medio millón de trabajadores fuera de la nómina estatal en seis meses, y de más de
un millón con el tiempo.

Los instrumentos eran reales. El trabajo por cuenta propia se amplió a 178 y luego
a 201 categorías, y el número de licencias pasó de unas 157.000 en 2010 a más de
medio millón para 2015 ---aproximadamente la quinta parte de la fuerza laboral
ocupada--- \citep{ritter2015}.\unv{} Se entregó tierra en usufructo bajo el
Decreto-Ley 259 de 2008, al principio 13 hectáreas, y bajo el Decreto-Ley 300 de
2012 hasta 67, a quien estuviera dispuesto a cultivarla; se repartieron algo más
de 1,7 millones de hectáreas ociosas.\unv{} Se autorizaron cooperativas no
agropecuarias a título de experimento en 2013. La reducción de plantilla nunca se
ejecutó a nada parecido a la velocidad anunciada, porque no había adónde mandar a
los trabajadores.

Las reformas se estancaron, y la razón por la que se estancaron es lo más útil de
este capítulo para la Parte III, así que conviene enunciarla con precisión en vez
de atribuirla a una renuencia general.

Un negocio privado cubano entre 2011 y 2021 no tenía personalidad jurídica. Era un
individuo con licencia, no una empresa. No podía por tanto firmar un contrato como
entidad, no podía demandar ni ser demandado, no podía poseer activos separados de
los de su dueño, y no podía, en la mayoría de los casos, venderse ni heredarse como
negocio en marcha.

No tenía mercado mayorista. Un restaurante compraba sus ingredientes en las mismas
tiendas minoristas que los hogares, a precios minoristas, compitiendo con los
hogares por los mismos bienes escasos ---lo que es a la vez económicamente absurdo
y el origen de la queja popular persistente de que el sector privado causó la
escasez---.

No tenía derecho a importar. Todo había que comprarlo en el país o traerlo en el
equipaje personal por las mulas, viajeros que hacían viajes repetidos con maletas,
que es como llegó durante una década una parte sustancial de los insumos del
sector privado cubano.

No tenía crédito. El préstamo bancario al sector privado se anunció, se rodeó de
condiciones y nunca se materializó a escala.

No tenía acceso a divisa a ninguna tasa utilizable.

Y operaba por enumeración, no por derecho: una actividad era legal si figuraba en
una lista que el Estado podía acortar en cualquier momento, y la lista estaba
escrita en el vocabulario de los oficios ---barbero, forrador de botones, payaso
de fiestas--- con las profesiones excluidas. Un desarrollador de software no tenía
categoría. Un estudio de diseño no tenía categoría.

El resultado es que Cuba obtuvo el sector privado que sus reglas estaban diseñadas
para producir: varios cientos de miles de microoperadores en comida, transporte,
alojamiento y servicios personales, y esencialmente nada en manufactura,
tecnología ni en nada que requiriera una entidad, un contrato y una cadena de
suministro. El capítulo~\ref{ch:capital} toma este diagnóstico como punto de
partida, porque casi todas las restricciones de la lista son legales e internas,
podrían haberse retirado sin un dólar de inversión extranjera, y no se retiraron
durante una década.

\section{La apertura}

El 17 de diciembre de 2014 Raúl Castro y Barack Obama anunciaron simultáneamente,
desde La Habana y Washington, que sus países restablecerían relaciones
diplomáticas. El intercambio inmediato fue la liberación del contratista
norteamericano Alan Gross y de un agente de inteligencia en manos cubanas, a
cambio de los tres miembros restantes de los Cinco; detrás había dieciocho meses
de conversaciones secretas mediadas en Canadá y por el Vaticano
\citep{leogrande2014}.

Las embajadas reabrieron en julio de 2015. Obama visitó La Habana en marzo de 2016
---primer presidente en ejercicio desde Coolidge en 1928--- y pronunció un
discurso, transmitido en vivo por la televisión cubana, en el que dijo que el
embargo terminaría y se dirigió directamente a los cubanos sobre su propio futuro.
Se reanudaron los vuelos regulares; atracaron cruceros; Airbnb entró en el mercado
en 2015 y en dos años tenía decenas de miles de anuncios cubanos; los visitantes
norteamericanos pasaron de unos 91.000 en 2014 a más de 600.000 en 2017, además de
un flujo mucho mayor de cubanoamericanos.\unv{}

Para el sector privado cubano fueron los mejores dieciocho meses que ha tenido
jamás. La demanda fue abrumadoramente a la economía no estatal, porque es donde
los norteamericanos con licencias de pueblo a pueblo estaban obligados a gastar:
restaurantes privados, cuartos privados, taxis privados, guías privados. El dinero
entró en los hogares y no en el Estado, entró equipamiento con los parientes de
visita, y por primera vez desde 1959 hubo una clase media cubana visible,
creciente y de propiedad nacional que se ganaba la vida con extranjeros sin el
Estado de intermediario.

Eso es también, precisamente, por lo que se detuvo.

\section{El freno}

El VII Congreso del Partido de abril de 2016 fue donde se frenaron los
lineamientos. El tono de los documentos pasó de la expansión del sector no estatal
a su regulación, y la frase que circuló fue la advertencia contra la concentración
de la riqueza y la propiedad. En agosto de 2017 el Estado suspendió la emisión de
nuevas licencias en un conjunto de las categorías más exitosas ---restaurantes
privados, alquiler de habitaciones y otras--- a la espera de una revisión que duró
más de un año y terminó en un régimen más estricto: una licencia por persona,
topes al tamaño de los restaurantes e inspección más dura.

Dos cosas más ocurrieron en ese mismo lapso corto. Fidel Castro murió el 25 de
noviembre de 2016. Y el 8 de noviembre de 2016 fue electo Donald Trump, que en
junio de 2017 emitió un memorando de política que revirtió una parte sustancial de
la apertura ---restringiendo los viajes individuales de pueblo a pueblo,
prohibiendo transacciones con una lista de entidades vinculadas a los militares, y
comenzando la secuencia de medidas que retoma el capítulo~\ref{ch:crisis}---.

Es tentador, y frecuente, atribuir el cierre enteramente a Washington. La
secuencia no lo respalda. La congelación de licencias de agosto de 2017 fue una
decisión cubana, tomada sobre negocios cubanos, antes de que mordieran las medidas
norteamericanas. La Habana y Washington llegaron a la misma conclusión por razones
opuestas: que el crecimiento de un sector privado cubano independiente tenía
consecuencias políticas. Uno lo quería y dejó de subvencionarlo; el otro no lo
quería y dejó de permitirlo.

Raúl Castro entregó la presidencia a Miguel Díaz-Canel en abril de 2018 y siguió
como Primer Secretario del partido hasta 2021. Dejó un país con un sector privado
mayor, una ciudadanía más libre y una relación funcional con Estados Unidos, y con
todos los problemas estructurales identificados en 2011 ---la dualidad monetaria,
la plantilla estatal, la autonomía empresarial, la producción de alimentos---
diagnosticados, anunciados y sin resolver.

\begin{ledger}{la década de Raúl, 2006--2018}
\emph{Logrado.} La derogación de prohibiciones cuyo único propósito era la
humillación: hoteles, teléfonos, computadoras, la venta de casas y autos y, sobre
todo, el permiso de salida, que después de cincuenta y cuatro años convirtió el
salir temporalmente de Cuba en un acto corriente. Un sector privado llevado de
150.000 a más de medio millón de personas, cerca de la quinta parte de la fuerza
laboral. Tierra ociosa repartida en usufructo a gran escala. Una consulta nacional
que la población se tomó en serio. Relaciones diplomáticas con Estados Unidos
restablecidas, embajadas reabiertas, y dieciocho meses en los que la demanda
norteamericana fluyó directamente hacia los hogares cubanos.

\emph{Costo.} Un programa de reformas cuyos propios instrumentos estaban diseñados
para que no pudiera tener éxito: sin personalidad jurídica, sin mercado mayorista,
sin derecho a importar, sin crédito, sin divisa, y con la legalidad por lista
enumerada. El resultado fueron varios cientos de miles de microoperadores y nada
por encima, que es exactamente lo que esas reglas producen. La única medida que
habría hecho legible la economía ---la unificación monetaria--- se anunció en 2011
y se aplazó una década, y acabaría intentándose en el peor momento posible. Y la
apertura que mejor funcionó se frenó desde dentro, en agosto de 2017, antes de que
Washington la cerrara desde fuera.
\end{ledger}
